\section{Proof of Theorem~\ref{thm:main}}
\label{sec:certificate}

We prove Theorem~\ref{thm:main} by applying
Proposition~\ref{fw:transfer} to the fields of
Theorem~\ref{tw:field-family}, with the upper bound for the relative
zeta value in Theorem~\ref{an:ceiling} and the data of
Definition~\ref{cert:data}. Propositions~\ref{cert:exact}
and~\ref{cert:enclosures} are verified by computer, in exact rational
arithmetic and in outward-rounded interval arithmetic, as described in
Subsection~\ref{cert:computations}; the other steps are proved in the
text. The inequalities between decimals in
Propositions~\ref{cert:exact} and~\ref{cert:enclosures} and in the remarks,
and the entries of Tables~\ref{cert:tab-local} and~\ref{cert:tab-budget},
are computed by the supplementary program
\texttt{margin\_general.py}, which calls the program \texttt{geom241.py} of
\cite{Naslund04273} unchanged; the replay records the final margin, and
calling \texttt{margin\_general} with \texttt{verbose=True} prints every
enclosure and table entry. In the inequalities the decimals are rounded
outward from interval enclosures, and the table entries are rounded to
ten decimals. No rounded decimal is substituted into a later calculation.

\subsection{Parameters and Local Data}

This subsection fixes $\delta$, $C$, the archimedean profiles and the shell
profiles (Definition~\ref{cert:data}), counts the places that carry shell
profiles (Lemma~\ref{cert:places}), and defines the computable lower bound
$\mathcal M_*$ for the margin (Definition~\ref{cert:lower-bounds}).

\begin{definition}\label{cert:data}
The \emph{data} of this section are the following.
\begin{enumerate}
\item[(a)] $\delta=43171/10^6=\deltastar$, $p=2/(1+\delta)=2000000/1043171$
 and $C=C_{\rm eff}=\Ceff$; $\lambda=2^{9/4}\sqrt{3615}$ and $\ell=\log\lambda$,
 as in Theorem~\ref{tw:field-family}.
\item[(b)] At the compact coordinates, $f_\R$ is the Gaussian of
 Lemma~\ref{pf:gaussian}, with $f_\R^p=e^{-a_\R|z|^2}$ and
 $a_\R=2p\delta=172684/1043171<1$.
\item[(c)] At the pair coordinates, $g_\C$ is the Student--Bernstein
 profile of Definition~\ref{pf:student} with $s=1.147487$,
 $a=0.0000181980309901$ and the polynomial $P$ of degree $m=3$ in each
 variable whose Bernstein coefficients are \eqref{pf:bernstein-witness}. We
 put $q_\star=sp-1=1251803/1043171$, as in Proposition~\fref{pf:mass}.
\item[(d)] In Definition~\fref{pf:tube-constant} and Lemma~\fref{pf:tube} we take
 $\rho=10^{-3}$, with the Bernstein coefficients of degree $m=3$ of (c). In
 Proposition~\fref{pf:mass} we partition $[0,1]^2$ into the $64$ squares of
 side $1/8$, use on each square the Bernstein coefficients of degree $3$, with
 $m_{\mathcal Q}$ and $M_{\mathcal Q}$ the smallest and largest of these
 coefficients, and take $N=18$. In Proposition~\fref{pf:overlap} we take
 $n_*=1$.
\item[(e)] For $K\in\mathcal K$, $\mathcal P$ is the set of places of $F$
 above the selected primes of $B$ (Definition~\ref{an:selected-def}). Let
 $\mathcal R=\{2,3,5,29,41\}$. We index the local data by $r\in\mathcal R$;
 the index $41$ always refers to the places above the capped prime
 $\mathfrak t_3$ of $B$, the only selected prime above $41$, and never to the
 places above $\mathfrak t_4$. For $r\in\mathcal R$ let $k_r$ be the
 exponent of Table~\ref{cert:tab-local}. At a place of $\mathcal P$ above
 $r\in\{2,3,5\}$, $g_r$ is the shell profile with parameters
 $(Q_r,k_r,(w_{r,ij}))$, with the symmetric weight matrices of
 Appendix~\ref{fw-data:section}; at a place above $29$ or above
 $\mathfrak t_3$ it is the shell profile with parameters
 $(Q_r,k_r,(x_{r,i}x_{r,j}))$, with the vectors $x_{29}$, $x_{41}$ of that
 appendix. Here $Q_r=r^4$ for $r=2,29,41$ and $Q_r=r^2$ for $r=3,5$, the
 order of the residue field of these places (Theorem~\ref{tw:field-family}).
 By Lemma~\ref{fw:shell-lemma}, $\mathcal F_{\delta,r}=\mathcal F_\delta(g_r)$
 depends only on these parameters.
\end{enumerate}
\end{definition}

Every parameter written as a finite decimal is exact. We recall
$\theta_*=\thetastar$ from Theorem~\ref{tw:field-family}.

\begin{definition}\label{cert:effective-types}
For the places of $\mathcal P$ we put $(e_r,f_r)=(8,4)$, $(2,2)$, $(2,2)$,
$(1,4)$ and $(2,4)$ for $r=2,3,5,29,41$. For $r\le29$ this is the absolute
type of the places of $\mathcal P$ above $r$. We call $(e_{41},f_{41})=(2,4)$
the \emph{effective type at $41$}: it is a counting device, not a local type.
The places of $\mathcal P$ above $41$ are the places above $\mathfrak t_3$;
they have absolute type $(1,4)$, and the effective type is chosen so that
their number is $d/(e_{41}f_{41})$ and their residue fields have order
$41^{f_{41}}$ (Lemma~\ref{cert:places}).
\end{definition}

For $r\in\{2,3,5,29\}$ these are uniform local types in the sense of
Definition~\ref{geo:uniform-types} (Theorem~\ref{tw:field-family}); because of
the index $41$, the proof of Theorem~\ref{thm:main} applies
Proposition~\fref{fw:transfer} directly rather than
Corollary~\fref{fw:uniform-types}.

\begin{lemma}\label{cert:places}
Let $K\in\mathcal K$, let $d=[F:\Q]$, let $\mathcal P$, $\mathcal R$, $Q_r$,
$k_r$ and $g_r$ be as in Definition~\ref{cert:data}(e), and let $(e_r,f_r)$ be
as in Definition~\ref{cert:effective-types}. For a place $\mathfrak p$ of
$\mathcal P$ above $r$ write $g_{\mathfrak p}=g_r$, $k_{\mathfrak p}=k_r$ and
$Q_{\mathfrak p}=Q_r$. Every place of $\mathcal P$ splits
in $K$. For $r\in\mathcal R$, the number of places of $\mathcal P$ above $r$
is $d/(e_rf_r)$, and each has residue field of order $Q_r$. Consequently
\[
 \frac1d\sum_{\mathfrak p\in\mathcal P}\log\mathcal F_\delta(g_{\mathfrak p})
 =\sum_{r\in\mathcal R}\frac{\log\mathcal F_{\delta,r}}{e_rf_r},\qquad
 \frac1d\sum_{\mathfrak p\in\mathcal P}k_{\mathfrak p}\log Q_{\mathfrak p}
 =\sum_{r\in\mathcal R}\frac{k_r\log r}{e_r}=:H.
\]
\end{lemma}

\begin{proof}
By Theorem~\ref{tw:field-family}, the places of $F$ above a selected prime
$\mathfrak p$ of $B$ split in $K$, and $\mathfrak p$ has a type $(e,f)$ in
$K/B$; so $K$ has $[K:B]/(ef)=d/(ef)$ primes above $\mathfrak p$, and $F$ has
half as many, each with residue field of order $\mathrm N\mathfrak p^f$. The
types are $(8,4)$ at the two primes above $2$, $(2,2)$ at the two primes above
$3$ and at the two above $5$, and $(1,4)$ at $\mathfrak t_1,\mathfrak t_2$
above $29$ and at $\mathfrak t_3$ above $41$. This gives $2\cdot d/64=d/32$
places above $2$, $2\cdot d/8=d/4$ above $3$, above $5$ and above $29$, and
$d/8$ above $41$ (the places above $\mathfrak t_3$, each with residue field
of order $41^4$), which is $d/(e_rf_r)$ in each case. The places above the
uncapped prime $\mathfrak t_4$, whose residue degree is a multiple of $4$
that may be larger, carry no shell profiles. The
second identity follows in the same way, since
$(d/(e_rf_r))\,k_r\log r^{f_r}/d=k_r\log r/e_r$.
\end{proof}

Since the weights are nonnegative with $w_{r,00}=1$ (Appendix~\ref{fw-data:section}
and Proposition~\ref{cert:exact}(a)), Lemma~\ref{fw:shell-lemma} gives
\begin{equation}\label{cert:finite-factor}
 \log\mathcal F_{\delta,r}
 =\log\Bigl(Q_r^{-\delta k_r}\sum_{(i,j),(i',j')}w_{r,ij}\mathbf G_{(i,j),(i',j')}w_{r,i'j'}\Bigr)
 -(1+\delta)\log\sum_{i,j}m_im_jw_{r,ij}^p,
\end{equation}
with the shell measures $m_i$ and the shell Gram matrix $\mathbf G$ of
Definition~\ref{fw:gram-matrix} for $Q=Q_r$ and $k=k_r$; for product weights this is
\eqref{fw:product-shells}. Table~\ref{cert:tab-local} lists the local data.

\begin{table}[htbp]
\centering\small
\caption{Local data. The last two columns give
$\log\mathcal F_{\delta,r}/(e_rf_r)$ for the shell profile $g_r$ and for
the unweighted shell profile with parameters $(Q_r,k_r)$, rounded to ten
decimals. The row $41$ refers to the places above $\mathfrak t_3$, with the
effective type $(e_{41},f_{41})$ of Definition~\ref{cert:effective-types}.}
\label{cert:tab-local}
\begin{tabular}{rrrrrlrr}
\toprule
$r$ & $e_r$ & $f_r$ & $Q_r$ & $k_r$ & weights & $g_r$ & unweighted\\
\midrule
$2$ & $8$ & $4$ & $16$ & $7$ & $10\times10$ matrix & $0.0390649879$ & $0.0387991734$\\
$3$ & $2$ & $2$ & $9$ & $9$ & $10\times10$ matrix & $0.3652797033$ & $0.3622194132$\\
$5$ & $2$ & $2$ & $25$ & $6$ & $10\times10$ matrix & $0.2795380841$ & $0.2780344049$\\
$29$ & $1$ & $4$ & $707281$ & $1$ & product, $3$ shells & $0.0279173524$ & $0.0279172669$\\
$41$ & $2$ & $4$ & $2825761$ & $1$ & product, $2$ shells & $0.0064840963$ & $0.0064840877$\\
\bottomrule
\end{tabular}
\end{table}

\begin{definition}\label{cert:lower-bounds}
Let $Z_*$ be the right side of \eqref{pf:collected-overlap} for the data
of Definition~\ref{cert:data}, with $n_*=1$. Let $\bar A_\C^*$ be the
upper end of the rigorous interval enclosure of $\bar A_\C$ computed by
the program of Subsection~\ref{cert:computations}
(Proposition~\ref{cert:enclosures}(a)), so that
$\bar A_\C\le\bar A_\C^*<45.4055724808268589$, the last number being
$\bar A_\C^*$ rounded up. Let $J_\C^*$ be the right side of
\eqref{pf:collected-functional} for these $Z_*$ and $\bar A_\C^*$, and put
\[
 \begin{split}
 \mathcal M_*(\theta)={}&\sum_{r\in\mathcal R}
 \frac{\log\mathcal F_{\delta,r}}{e_rf_r}
 -(1/2-\delta)\ell-C+(1-\delta)\log2+(1-\theta)\log\pi\\
 &+(1-2\theta)J_\R+\theta J_\C^* .
 \end{split}
\]
\end{definition}

\subsection{Finite Checks and Rigorous Bounds}

This subsection states the finite checks (Proposition~\ref{cert:exact}) and
the interval enclosures (Proposition~\ref{cert:enclosures}) on which the proof
rests; both are verified by computer as described in
Subsection~\ref{cert:computations}.

\begin{proposition}\label{cert:exact}
For the data of Definition~\ref{cert:data}:
\begin{enumerate}
\item[(a)] all weights $w_{r,ij}$ ($r=2,3,5$, $0\le i,j\le9$) and
 $x_{r,i}$ ($r=29,41$) are positive (Lemma~\fref{fw:shell-lemma} needs only
 $w_{r,ij}\ge0$ with $w_{r,00}>0$, which is what the programs assert; strict
 positivity is checked separately);
\item[(b)] the coefficients \eqref{pf:bernstein-witness} are positive,
 and the tube constant satisfies $q_0<0.0296$;
\item[(c)] on each of the $64$ squares $\mathcal Q$, the Bernstein
 coefficients of degree $3$ of the restriction of $P$ to $\mathcal Q$ are
 positive, and $\rho_{\mathcal Q}<0.3484$;
\item[(d)] every cross satisfies \eqref{pf:cross-threshold}.
\end{enumerate}
\end{proposition}

\begin{proof}
These are finite checks, made by the supplementary programs
\texttt{margin\_general.py} and \texttt{geom241.py} as described in
Subsection~\ref{cert:computations}. The comparisons in (a)--(c), and the
first condition in (d), are exact comparisons of rational numbers; the other
two conditions in (d) are checked by outward-rounded interval evaluation of
$\log(4/a)$.
\end{proof}

\begin{proposition}\label{cert:enclosures}
For the data of Definition~\ref{cert:data}, with $Z_*$, $J_\C^*$ and
$\mathcal M_*$ as in Definition~\ref{cert:lower-bounds}, the following
hold.
\begin{enumerate}
\item[(a)] With $T$ and $T'$ independent with density $q_\star t^{q_\star-1}$ on
$(0,1)$,
\begin{equation}\label{cert:mass}
 45.4055724808268557<\bar A_\C=\mathbb E\,P(T,T')^p<45.4055724808268589 ,
\end{equation}
and the right side of \eqref{pf:subrectangle-mass} is less than
$1.6\cdot10^{-15}$.
\item[(b)] The expression $Z_*$ satisfies
\begin{equation}\label{cert:overlap}
 397.6858754715135886<Z_*<397.6858754715135887 ,
\end{equation}
and the remainder term satisfies $E_{-,1}<2.8\cdot10^{-16}$.
\item[(c)] $1.3444924010459460<J_\C^*<1.3444924010459462$.
\item[(d)] $-0.2125007533361271<J_\R<-0.2125007533361270$.
\item[(e)] $0.7182842240056688517<\sum_{r\in\mathcal R}\log\mathcal
F_{\delta,r}/(e_rf_r)<0.7182842240056688518$.
\item[(f)] The numbers $H$ of Lemma~\ref{cert:places} and
$\ell=\frac94\log2+\frac12\log3615$ satisfy
\begin{gather*}
 15.6026546826373748<H<15.6026546826373749,\\
 5.6560047235563094<\ell<5.6560047235563095 .
\end{gather*}
\item[(g)] $\log K_\C<0.0558518577$, $\varsigma_\C>1.4728810283$ and
$\mu=2e^{H/2-\ell}>17.0895949922$.
\item[(h)] $0.6247640218687999<-\log\pi-2J_\R+J_\C^*<0.6247640218688001$.
\item[(i)] The lower bound for the margin satisfies
\begin{equation}\label{cert:margin}
 \mathcal M_*(\theta_*)>0.0000105426526168 .
\end{equation}
\end{enumerate}
\end{proposition}

\begin{proof}
The program \texttt{margin\_general.py} evaluates these quantities in
outward-rounded interval arithmetic from the exact data, as described in
Subsection~\ref{cert:computations}: (a) by Proposition~\ref{pf:mass}; (b)
from \eqref{pf:collected-overlap}, with Lemma~\ref{pf:digamma} for the
digamma values; (c) from \eqref{pf:collected-functional}; (d) from
\eqref{pf:gaussian-compact}; (e) from \eqref{cert:finite-factor}; (f) from
the definitions of $H$ and $\ell$; (g) from Lemma~\ref{pf:tube} and the
definition of $\mu$; and (h) and (i) from
Definition~\ref{cert:lower-bounds}, where $\bar A_\C^*$ is the upper end of
the enclosure computed for (a).
\end{proof}

The proof of Theorem~\ref{thm:main} uses Proposition~\ref{cert:exact} and,
from Proposition~\ref{cert:enclosures}, only the fact, established by the
computation for (a), that $\bar A_\C\le\bar A_\C^*$, the positivity of $Z_*$
given by (b), and (g), (h) and (i). Table~\ref{cert:tab-budget} lists the
terms of $\mathcal M_*(\theta_*)$.

\begin{table}[htbp]
\centering\small
\caption{The terms of $\mathcal M_*(\theta_*)$, rounded to ten decimals.}
\label{cert:tab-budget}
\begin{tabular}{lr}
\toprule
Term & Value\\
\midrule
$\sum_r\log\mathcal F_{\delta,r}/(e_rf_r)$ & $0.7182842240$\\
$-(1/2-\delta)\ell$ & $-2.5838269819$\\
$-C$ & $-0.0422764000$\\
$(1-\delta)\log2$ & $0.6632233236$\\
$(1-\theta_*)\log\pi$ & $0.5723736765$\\
$(1-2\theta_*)J_\R$ & $-0.0000032425$\\
$\theta_*J_\C^*$ & $0.6722359429$\\
\midrule
$\mathcal M_*(\theta_*)$ & $0.0000105427$\\
\bottomrule
\end{tabular}
\end{table}

\subsection{Conclusion of the Proof}

This subsection verifies the hypotheses of Proposition~\ref{fw:transfer} for
the fields of Theorem~\ref{tw:field-family} and the data of
Definition~\ref{cert:data}, and concludes the proof of
Theorem~\ref{thm:main}.

\begin{corollary}\label{cert:fourier}
The pair $(f_\R,g_\C)$ of Definition~\ref{cert:data} is admissible
(Definition~\ref{geo:admissible}) with
Fourier constants $M=2$ and $\varsigma=1$, and these satisfy
\eqref{fw:fourier-condition} with $H_f=H$.
\end{corollary}

\begin{proof}
Proposition~\ref{pf:admissible} applies, since $s\ge1$, $a>0$,
$a_\R\le1$ and $0<\rho<1$, the Bernstein coefficients
\eqref{pf:bernstein-witness} of degree $m=3$ are positive, and $q_0<1$
(Proposition~\ref{cert:exact}(b)). It gives the Fourier constants
$M=\max\{2,\sqrt{K_\C}\}$ and $\varsigma=\min\{1,\varsigma_\C\}$. By
Proposition~\ref{cert:enclosures}(g), $K_\C<4$ and $\varsigma_\C>1$, so
$M=2$ and $\varsigma=1$. The right side of \eqref{fw:fourier-condition} is
then $\log2+2\log5+2<5.9121$, while
$\varsigma\cdot2e^{H_f/2-\ell}=\mu>17.0895949922$ by Lemma~\ref{cert:places}
and Proposition~\ref{cert:enclosures}(g).
\end{proof}

\begin{lemma}\label{cert:lower-margin}
For the data of Definition~\ref{cert:data}, $J_\C\ge J_\C^*$.
Consequently, for $\theta_*\le\theta<1/2$, the expression
$\mathcal M_{\mathcal R}(\theta)$ obtained from $\mathcal M_*(\theta)$ by
replacing $J_\C^*$ with $J_\C$ (by Lemma~\ref{cert:places}, at $\theta=c/d$
for $F$ it is the margin \eqref{fw:margin} of $K/F$, for $K\in\mathcal K$;
the same expression appears in Corollary~\fref{fw:uniform-types}) satisfies
\[
 \mathcal M_{\mathcal R}(\theta)\ge\mathcal M_*(\theta)\ge\mathcal
 M_*(\theta_*)>0.0000105426526168 .
\]
\end{lemma}

\begin{proof}
Corollary~\ref{pf:functional} applies: the hypotheses of
Propositions~\ref{pf:overlap} and~\ref{pf:mass} hold by
Proposition~\ref{cert:exact}(c),(d), $Z_*>0$ by \eqref{cert:overlap},
and $\bar A_\C\le\bar A_\C^*$ by Proposition~\ref{cert:enclosures}(a) and
Definition~\ref{cert:lower-bounds}. Hence $J_\C\ge J_\C^*$. Since
$\mathcal M_{\mathcal R}(\theta)-\mathcal M_*(\theta)=\theta(J_\C-J_\C^*)$
and $\theta\ge0$, we get $\mathcal M_{\mathcal R}(\theta)\ge\mathcal
M_*(\theta)$. The function $\mathcal M_*$ is affine in $\theta$ with slope
$-\log\pi-2J_\R+J_\C^*$, which is positive by
Proposition~\ref{cert:enclosures}(h); so $\mathcal M_*(\theta)\ge\mathcal
M_*(\theta_*)$ for $\theta\ge\theta_*$, and \eqref{cert:margin} completes the
proof.
\end{proof}

\begin{proof}[Proof of Theorem~\ref{thm:main}]
For each $n\ge49$, let $K^{(n)}/B$ be a Galois extension with
$[K^{(n)}:B]=2^n$ as in Theorem~\ref{tw:field-family}, let $\iota_1$ be
a complex conjugation at a place of $K^{(n)}$ above the real place $v_1$
of $B$, put $F^{(n)}=(K^{(n)})^{\langle\iota_1\rangle}$ and
$d_n=[F^{(n)}:\Q]$, and let $\theta_n$ be the value of $\theta$ for
$F^{(n)}$. By Theorem~\ref{an:ceiling}, $d_n^{-1}\log L_{F^{(n)}}(1)<C$
for all sufficiently large $n$; we discard the finitely many other
fields.

We apply Proposition~\ref{fw:transfer} to this sequence, with
$\delta=\deltastar$, with $\mathcal P_n$ the set $\mathcal P$ of
Definition~\ref{cert:data}(e) for $K^{(n)}$, the shell profiles $g_r$ at its
places, and the pair $(f_\R,g_\C)$ of Definition~\ref{cert:data}, and check
its hypotheses. By Theorem~\ref{tw:field-family}, each $K^{(n)}/F^{(n)}$
satisfies (G1)--(G3) with the common $\lambda=2^{9/4}\sqrt{3615}$, and
$d_n=2^n\to\infty$. By the choice above, $\log L_{F^{(n)}}(1)\le d_nC$. The
places of $\mathcal P_n$ split in $K^{(n)}$ (Lemma~\ref{cert:places}). The
weights are as in Definition~\ref{fw:shell-profile}, by
Proposition~\ref{cert:exact}(a), and the parameters of the shell profiles take
five values, independent of $n$. The pair $(f_\R,g_\C)$ is admissible, and its
Fourier constants satisfy \eqref{fw:fourier-condition}, by
Corollary~\ref{cert:fourier}, since $H_f=H$ for every $n$ by
Lemma~\ref{cert:places}. By Lemma~\ref{cert:places}, the margin
\eqref{fw:margin} of $K^{(n)}/F^{(n)}$ with these shell profiles is
$\mathcal M_{\mathcal R}(\theta_n)$. Finally, $\theta_*\le\theta_n<1/2$ by
Theorem~\ref{tw:field-family}, so Lemma~\ref{cert:lower-margin} gives
$\mathcal M_{\mathcal R}(\theta_n)\ge\mathcal M_*(\theta_*)>0$ for all $n$.

Proposition~\ref{fw:transfer} therefore gives finite sets
$U_n\subset\R^2$ with $|U_n|\to\infty$ and
$u(U_n)/|U_n|^{1+\delta}\to\infty$, where $1+\delta=1.043171$; these are
the sets of Theorem~\ref{thm:main}, indexed by $n$. In particular
$u(U_n)\ge|U_n|^{1.043171}$ for all large $n$, which gives the last
assertion of Theorem~\ref{thm:main} at the cardinalities $|U_n|$.
\end{proof}

\subsection{Remarks}

This subsection records numerical observations that are not used in the
proof.

\begin{remark}\label{cert:signature-remark}
The program also encloses the value at $\theta=0$ and the zero $\theta_0$
of the affine function $\mathcal M_*$:
\[
 \mathcal M_*(0)<-0.312,\qquad \theta_0<0.49997550 .
\]
These enclosures are not used in the proof. Since $\theta=1/2-1/(4N_\iota)$
by Theorem~\ref{tw:field-family}, where $N_\iota$ is the number of conjugates
of $\iota_1$ in $\Gal(K/B)$, $N_\iota\ge10205$ would suffice, whereas that
theorem gives $N_\iota\ge2^{15}$. At $\theta=0$ the margin does not involve
$J_\C$ and is negative: with these profiles, mixed signature is essential.
\end{remark}

\begin{remark}\label{cert:observations}
The following numerical observations are not used in the proof. The proof of
Proposition~\ref{fw:transfer} needs only some $\eta>0$ with $4\eta$ below the
margin; for instance, $\eta=10^{-9}$ gives
\begin{equation}\label{cert:final-margin}
 \mathcal M_*(\theta_*)-4\eta>0.0000105386526168>0 .
\end{equation}
The value $\mathcal M_*(\theta_*)$ would be absorbed by an increase of
$1.05\cdot10^{-5}$ in $C$. The general weights at $2$, $3$ and $5$ add
$1.53\cdot10^{-5}$ to the margin: with the product weights of the
intermediate certificate, optimized at $\delta=0.043170$, the margin at
$\delta=\deltastar$ is about $-4.8\cdot10^{-6}$. At $\delta=0.043172$, with the
same data and the same $C$, the evaluation gives about $-1.33\cdot10^{-5}$.
\end{remark}

\subsection{Computations}\label{cert:computations}
This subsection describes how the finite checks and enclosures of
Propositions~\ref{cert:exact} and~\ref{cert:enclosures} are computed.

The program \texttt{margin\_general.py} calls the function \texttt{margin} of
\texttt{geom241.py} of \cite{Naslund04273} unchanged, and replaces only, for
the duration of the call, its routine for the local functional at the places
whose weights are given as a matrix, by the routine of
\texttt{general\_windows.py}, which evaluates \eqref{fw:shell-kernel} exactly
in rational arithmetic, with the real powers $w_{r,ij}^p$ in interval
arithmetic. The routine for product weights, and the routines for
Propositions~\ref{pf:mass} and~\ref{pf:overlap}, Lemma~\ref{pf:tube} and
Corollary~\ref{pf:functional}, are those of \cite{Naslund04273}, which take
three modules of the supplementary archive of \cite{Naslund0418235} unchanged.
In Proposition~\ref{pf:mass} the program uses the symmetry of
\eqref{pf:bernstein-witness}: each square off the diagonal is evaluated once
and counted twice. It checks the hypotheses listed in
Proposition~\ref{cert:exact} (for (c), on the squares on or above the
diagonal; the others follow by symmetry), the conditions $a_\R\le1$,
$\varsigma_\C>1$, $\log K_\C<2\log2$ and $\mu>10$ (which implies
\eqref{fw:fourier-condition} for $M=2$ and $\varsigma=1$, since $\log50+2<10$),
and the positivity of the slope in Proposition~\ref{cert:enclosures}(h). It
then prints the enclosures stated in this section. The exponents and weights
were produced by separate optimization programs; they enter the proof only
as the exact data of Appendix~\ref{fw-data:section}.

Rational quantities are computed exactly; the remaining quantities involve the
logarithm, the exponential (also through the real powers $x^y=\exp(y\log x)$),
the square root in $\varsigma_\C$, the constant $\pi$ and the digamma function,
which enters only through Lemma~\ref{pf:digamma}. These are evaluated in the
interval arithmetic of mpmath \cite{Mpmath}, with every operation rounded
outward, at a working precision of $80$ decimal digits. Hence every computed
interval contains the exact value, and the decimals stated above are rigorous
bounds.

The replay program \texttt{dihedral/reproduce\_w3.py} recomputes
Propositions~\ref{cert:exact} and~\ref{cert:enclosures} from the witness file,
and also runs \texttt{general\_windows\_check.py}
(Remark~\ref{fw:general-weights}).
